% =====================================================================
\BookChapter[How Code Becomes Computation]{How Code Becomes Computation}
\label{ch:howcoderuns}

\begin{researchquestion}
What happens between pressing Run and seeing a number, and why is the number sometimes slightly wrong?
\end{researchquestion}
% =====================================================================

When you press ``Run'' in a \py editor, a lot happens before your answer appears. This
chapter opens the box. We look at how source code is translated and executed, how a
computer stores numbers and text, and why some calculations give surprising answers.
These ideas are the same for every language, and they explain many of the differences
you will see in later chapters.

\LearningOutcomes{Explain how compilers, interpreters and just-in-time compilers execute a program; state the range of an integer type and the precision of a floating-point type; predict when rounding error and cancellation will damage a result; explain how arrays are laid out in memory and why loop order matters; distinguish syntax, run-time and logic errors; and name the libraries on which numerical software is built.}

\section{The computer in one page}

For our purposes a computer has four parts.
\begin{itemize}
\item The \term{processor} (CPU) executes machine instructions. A modern laptop CPU has
  several \emph{cores}, each of which can run its own stream of instructions, and each
  core performs billions of simple operations per second.
\item The \term{memory} (RAM) stores the program and its data while it runs. It is fast
  but forgets everything when the power goes off. A laptop typically has 8--32~GB.
\item The \term{storage} (SSD or hard disk) keeps files permanently. It is much slower
  than RAM.
\item \term[input/output devices]{Input/output devices}: keyboard, screen, network, sensors, and so on.
\end{itemize}
Many computers also have a \term{GPU} (graphics processing unit), which has thousands of
small cores that perform the same operation on many numbers at once. GPUs were built for
graphics but are now essential for machine learning and for some simulations.

Everything in memory is stored as \term{bits} (0 or 1), grouped into \term{bytes} of 8
bits. One byte can hold $2^8 = 256$ different values. Whether a particular group of bytes
means a number, a letter, a pixel colour or a machine instruction depends entirely on how
the program interprets it. That is what \emph{types} are for.

\section{Compilers}

A \term{compiler} reads the complete source code and produces an executable file of
machine code. Figure~\ref{fig:compile} shows the steps for a C program; Fortran, C++,
Rust and Go work the same way.

\begin{figure}[htbp]
\centering
\begin{tikzpicture}[node distance=7mm and 9mm,
  box/.style={draw=nublue,fill=nublue!6,rounded corners=2pt,minimum height=9mm,
    font=\sffamily\small,align=center,inner sep=4pt},
  file/.style={draw=gray,fill=white,font=\ttfamily\small,minimum height=7mm,inner sep=3pt},
  arr/.style={-{Stealth[length=2.5mm]},thick}]
\node[file] (src) {drag.c};
\node[box,right=of src] (pre) {pre-\\processor};
\node[box,right=of pre] (comp) {compiler};
\node[file,right=of comp] (obj) {drag.o};
\node[box,right=of obj] (link) {linker};
\node[file,right=of link] (exe) {drag};
\node[file,below=of link] (lib) {libm.so};
\draw[arr] (src) -- (pre);
\draw[arr] (pre) -- (comp);
\draw[arr] (comp) -- (obj);
\draw[arr] (obj) -- (link);
\draw[arr] (link) -- (exe);
\draw[arr] (lib) -- (link);
\node[below=1mm of comp,font=\sffamily\scriptsize,align=center,text=gray]{checks syntax and types,\\optimises, emits\\machine code};
\node[below=1mm of exe,font=\sffamily\scriptsize,align=center,text=gray]{run it:\\\texttt{./drag}};
\end{tikzpicture}
\caption{From source code to executable. The \emph{compiler} turns each source file into
an object file of machine code; the \emph{linker} joins object files and libraries (here
the maths library) into one executable.}
\label{fig:compile}
\end{figure}

On the command line it looks like this:
\begin{lstlisting}[style=shell]
$ gcc -O2 -o drag drag.c -lm      # compile and link, with optimisation
$ ./drag                          # run the executable
\end{lstlisting}
The option \code{-O2} asks the compiler to \term{optimise}: to rearrange and simplify the
machine code so that it runs faster while giving the same results. Optimising compilers
are remarkably good; they unroll loops, keep values in registers, and use special
instructions that process several numbers at once (called \term{SIMD} or vector
instructions).

Compiling has three consequences that you will notice when you use a compiled language.
\begin{enumerate}
\item \textbf{Errors are found early.} Spelling mistakes, missing semicolons and type
  mismatches stop the compilation; the program never runs with them.
\item \textbf{There is an extra step.} After every change you must recompile. For large
  programs this can take minutes, so tools such as \code{make} and \code{CMake} recompile
  only what has changed.
\item \textbf{The executable is tied to a platform.} A program compiled on Linux for an
  Intel processor does not run on a Mac with an Apple-silicon processor. You must compile
  it again there (the source code is portable; the executable is not).
\end{enumerate}

\section{Interpreters}

An \term{interpreter} executes a program directly from its source code, statement by
statement. The standard \py implementation, called \term{CPython}, first translates your
code into a compact \term{bytecode} and then runs a loop that reads one bytecode
instruction at a time and performs it. You can see the bytecode with the \code{dis}
module:
\begin{lstlisting}[language=Python]
import dis
def f(a, b):
    return a * b + 1
dis.dis(f)
\end{lstlisting}
The output (slightly shortened) lists instructions such as \code{LOAD\_FAST a},
\code{LOAD\_FAST b}, \code{BINARY\_OP *}, \code{LOAD\_CONST 1}, \code{BINARY\_OP +},
\code{RETURN\_VALUE}. Each of these is carried out by a piece of C code inside the
interpreter.

Interpreters have the opposite strengths of compilers.
\begin{itemize}
\item \textbf{Fast feedback}: you change a line and run it immediately. You can even type
  statements one at a time in an interactive session (a \term{REPL}: read--evaluate--print
  loop) or a Jupyter notebook cell.
\item \textbf{Portability}: the same \file{.py} file runs on Windows, macOS, Linux and a
  Raspberry Pi, as long as a \py interpreter is installed.
\item \textbf{Slower execution of loops}: every operation goes through the interpreter
  loop, checks the types of its operands, and creates new objects. A simple numerical loop
  in pure \py can be 10 to 100 times slower than the same loop in C.
\end{itemize}

\section{Virtual machines and just-in-time compilation}

Java and C\# take a middle path. Their compilers translate source code into bytecode for a
\term{virtual machine}: the Java Virtual Machine (JVM) or the .NET Common Language Runtime
(CLR). The bytecode is the same on every platform. When the program runs, the virtual
machine watches which parts are executed often and compiles those parts to native machine
code on the fly. This is \term{just-in-time (JIT) compilation}. After a short warm-up,
JIT-compiled code runs nearly as fast as C.

JavaScript engines in web browsers (V8 in Chrome and Node.js, SpiderMonkey in Firefox)
are JIT compilers too, which is why complex applications can run inside a web page.
\textbf{Julia} uses JIT compilation in a particular way: the first time you call a
function with arguments of certain types, Julia compiles a specialised machine-code version
for exactly those types. That first call is slow; later calls are fast.

\begin{table}[htbp]
\centering
\caption{How the languages in these notes are usually executed.}
\label{tab:execution}
\small
\begin{tabular}{@{}lll@{}}
\toprule
\textbf{Language} & \textbf{Usual execution model} & \textbf{File you run}\\
\midrule
Fortran, C, C++, Rust, Go & ahead-of-time compiled & native executable\\
Java, Kotlin, Scala & bytecode + JIT on the JVM & \file{.class} / \file{.jar}\\
C\#, F\#, Visual Basic .NET & bytecode + JIT on .NET & \file{.dll} / \file{.exe}\\
JavaScript, TypeScript & JIT in browser or Node.js & \file{.js} source\\
\py & bytecode interpreter (CPython) & \file{.py} source\\
MATLAB & interpreter with JIT engine & \file{.m} source\\
R, Wolfram Language & interpreter & \file{.R}, \file{.wl}/\file{.nb}\\
Julia & JIT (per-type specialisation) & \file{.jl} source\\
\bottomrule
\end{tabular}
\end{table}

\begin{keyidea}
The speed difference between languages matters mostly for \emph{loops over many simple
operations}. If your program spends its time inside a library routine (solving a linear
system, computing an FFT, training a neural network), it runs at compiled speed no matter
which language called the routine. This is why slow-looking languages like \py and
MATLAB are used for heavy numerical work: the heavy part is done by compiled Fortran, C
or C++ underneath.
\end{keyidea}

\section{How integers are stored}

An integer is stored in binary. With $n$ bits you can represent $2^n$ different values.
The number $13$ is $1101_2 = 1\cdot 8 + 1 \cdot 4 + 0 \cdot 2 + 1 \cdot 1$. Most languages
store integers in a fixed number of bits:
\begin{center}
\small
\begin{tabular}{@{}llr@{}}
\toprule
\textbf{Size} & \textbf{Typical names} & \textbf{Range (signed)}\\
\midrule
8 bits & \code{int8}, \code{signed char}, \code{byte} & $-128$ to $127$\\
16 bits & \code{int16}, \code{short} & $-32\,768$ to $32\,767$\\
32 bits & \code{int32}, \code{int} (C, Java), \code{integer} (Fortran) & about $\pm 2.1\times 10^{9}$\\
64 bits & \code{int64}, \code{long long}, \code{long} (Java) & about $\pm 9.2\times 10^{18}$\\
\bottomrule
\end{tabular}
\end{center}
Negative numbers are stored in \term{two's complement} form, which makes addition of
positive and negative numbers use the same circuit.

What happens if a result does not fit? This is called \term{integer overflow}. In C a
signed overflow is formally ``undefined behaviour'' --- anything may happen; in practice
the value usually wraps around:
\begin{lstlisting}[language=Python]
import numpy as np
print(np.int32(2147483647) + np.int32(1))     # -2147483648 (with a warning)
\end{lstlisting}
Plain \py integers are different: they grow automatically to as many bits as needed, so
\code{2**200} is computed exactly. This is convenient but slower, and it is one reason why
NumPy arrays use fixed-size integers.

\begin{history}[A little history: Ariane 5]
On 4 June 1996 the first Ariane 5 rocket veered off course and was destroyed 37 seconds
after launch. The inquiry found that a 64-bit floating-point number related to the
horizontal velocity was converted to a 16-bit signed integer. The value was larger than
$32\,767$, the conversion failed, and the guidance computers shut down. The code had been
reused from Ariane 4, whose flight profile never produced such large values. The rocket
and its payload were lost. A type conversion is an engineering decision.
\end{history}

\section{How real numbers are stored: floating point}
\label{sec:floatingpoint}

Real numbers such as $9.81$ or $6.022\times10^{23}$ are stored as \term{floating-point
numbers}, following the IEEE~754 standard that virtually all hardware uses. A
floating-point number is stored in scientific notation in base 2:
\[
  x = (-1)^s \times (1.f)_2 \times 2^{e},
\]
with a sign bit $s$, a fraction $f$ (the significant digits) and an exponent $e$. The two
common formats are
\begin{center}
\small
\begin{tabular}{@{}lllll@{}}
\toprule
\textbf{Format} & \textbf{Names} & \textbf{Bits} & \textbf{Significant digits} & \textbf{Largest value}\\
\midrule
single & \code{float}, \code{real(4)}, \code{float32} & 32 & about 7 & $3.4\times10^{38}$\\
double & \code{double}, \code{real(8)}, \code{float64} & 64 & about 16 & $1.8\times10^{308}$\\
\bottomrule
\end{tabular}
\end{center}
\py's \code{float}, MATLAB's default numbers, R's \code{numeric}, JavaScript's
\code{Number} and Julia's \code{Float64} are all IEEE double precision. In C, C++, Java and
C\# you choose \code{double} explicitly; in Fortran you choose a \code{kind}
(Section~\ref{sec:fortran-kinds}). For engineering calculations, use double precision
unless you have a good reason not to.

\subsection*{Most decimal numbers cannot be stored exactly}

Just as $1/3 = 0.3333\ldots$ cannot be written exactly with a finite number of decimal
digits, $0.1$ cannot be written exactly with a finite number of binary digits. The
computer stores the nearest representable value:
\begin{lstlisting}[language=Python]
>>> format(0.1, '.20f')
'0.10000000000000000555'
>>> 0.1 + 0.2
0.30000000000000004
>>> 0.1 + 0.2 == 0.3
False
\end{lstlisting}
This is not a bug in \py; you get the same answer in C, Java, JavaScript, MATLAB and every
other language that uses IEEE doubles. Small \term{rounding errors} like this can also
accumulate:
\begin{lstlisting}[language=Python]
total = 0.0
for i in range(10):
    total += 0.1
print(total, total == 1.0)        # 0.9999999999999999 False
\end{lstlisting}
(Interestingly, \py's built-in \code{sum([0.1]*10)} returns exactly \code{1.0} in
recent versions, because it uses a more careful summation algorithm internally. Libraries
care about these details so that you do not have to.)

\begin{pitfall}[Common pitfall: comparing floating-point numbers with \texttt{==}]
Never test whether two computed floating-point numbers are \emph{exactly} equal. Test
whether they are close, with a tolerance that makes sense for your problem:
\begin{lstlisting}[language=Python]
import math
math.isclose(0.1 + 0.2, 0.3)                 # True (relative tolerance 1e-9)
abs(a - b) <= 1e-12 * max(abs(a), abs(b))    # the same idea by hand
\end{lstlisting}
Every language has an equivalent: \code{numpy.isclose}, \code{all.equal} in R,
\code{isapprox} (or \code{a $\approx$ b}) in Julia, and \code{abs(a-b) < tol} in C,
Fortran and MATLAB.
\end{pitfall}

\subsection*{Machine epsilon sets the best possible precision}

The gap between $1$ and the next larger double is called \term{machine epsilon},
$\varepsilon = 2^{-52} \approx 2.22\times10^{-16}$. It sets the best relative precision you
can hope for. Adding a number smaller than about $\varepsilon x$ to $x$ has no effect:
\begin{lstlisting}[language=Python]
>>> import sys
>>> sys.float_info.epsilon
2.220446049250313e-16
>>> 1e16 + 1 - 1e16
0.0
\end{lstlisting}

\subsection*{Subtracting nearly equal numbers destroys digits}

The most dangerous rounding error happens when you subtract two nearly equal numbers. The
leading digits cancel, and what is left is dominated by rounding error. For small $x$,
$1-\cos x \approx x^2/2$, so $(1-\cos x)/x^2$ should be close to $0.5$:
\begin{lstlisting}[language=Python]
>>> import math
>>> x = 1e-8
>>> (1 - math.cos(x)) / x**2          # catastrophic cancellation
0.0
>>> 2 * math.sin(x/2)**2 / x**2       # same formula, rewritten
0.5
\end{lstlisting}
Since $\cos(10^{-8})$ is equal to $1$ to within machine precision, the first formula
loses every significant digit. The identity $1-\cos x = 2\sin^2(x/2)$ avoids the
subtraction. Numerical analysts spend careers on such rewrites; libraries such as
\code{math.expm1} and \code{math.log1p} exist for the same reason.

\subsection*{Infinity and not-a-number spread silently}

IEEE arithmetic does not stop on errors. Instead it produces special values:
\code{inf} (infinity, e.g.\ $1.0/0.0$ in C or NumPy), \code{-inf}, and \code{nan} (``not
a number'', e.g.\ $0.0/0.0$ or $\sqrt{-1}$). A \code{nan} spreads through every
calculation it touches, and \code{nan == nan} is \code{False}. If a simulation suddenly
prints \code{nan}, something earlier divided by zero or overflowed --- often because the
time step was too large.

\begin{engnote}
Single precision (\code{float32}) has only about 7 significant digits. That is enough
for a measured temperature, and GPUs process single (and even half) precision much faster
than double, which is why machine-learning code uses it. But when a simulation adds many
small increments to a large total (a long time integration, a large sum of forces), 7
digits can be far too few. If results change when you switch from \code{float32} to
\code{float64}, believe the \code{float64} results.
\end{engnote}

\section{Text: characters and Unicode}

Text is stored as numbers too. The old \term{ASCII} code assigned the numbers 0--127 to
the English letters, digits and punctuation; \code{\sq{}A\sq{}} is 65. To represent all the
world's writing systems --- Kazakh, Russian, Chinese, mathematical symbols, emoji ---
modern systems use \term{Unicode}, which assigns a number (a \emph{code point}) to more
than 150\,000 characters. The \term{UTF-8} encoding stores these code points in one to
four bytes and is the standard for files and the web. \py~3 strings are Unicode by
default, so strings that contain Kazakh or Russian letters, Greek symbols such as
$\Delta$ or $\mu$, or the degree sign work without any effort. In C
and older Fortran, text is a sequence of bytes, and handling non-English text needs more
care.

\section{Memory: names, values and references}

Different languages have different mental models of what a variable is.

\paragraph{Variables as boxes (C, Fortran).}
In C, the declaration \code{double x;} reserves a box of 8 bytes in memory, and the name
\code{x} refers to that box. \code{x = 3.0;} puts a value into the box. \code{y = x;}
\emph{copies} the value into another box.

\paragraph{Variables as labels (\py, JavaScript, Java objects).}
In \py, \code{x = [1, 2, 3]} creates a list object somewhere in memory and attaches the
name \code{x} to it. \code{y = x} attaches a second name to the \emph{same} object --- it
does not copy it:
\begin{lstlisting}[language=Python]
x = [1, 2, 3]
y = x
y.append(4)
print(x)          # [1, 2, 3, 4]  -- x and y are the same list
z = x.copy()      # an independent copy
\end{lstlisting}
This is called \term{reference semantics}. MATLAB and R behave differently: they use
\term{value semantics} for arrays, so \code{y = x; y(1) = 99;} leaves \code{x} unchanged
(internally they copy lazily, only when you modify). Knowing which model a language uses
prevents some very confusing bugs.

\paragraph{The stack and the heap.}
Programs use two regions of memory. The \term{stack} holds the local variables of
functions that are currently running; it is fast and is cleaned up automatically when a
function returns. The \term{heap} holds data whose size or lifetime is not known in
advance, such as a large array read from a file. In C you request heap memory with
\code{malloc} and must release it with \code{free}; forgetting to do so is a \term{memory
leak}, and using memory after freeing it is a serious bug. \py, Java, C\#, JavaScript,
Julia, R and MATLAB use a \term{garbage collector} that frees unused memory automatically.
C++ and Rust use rules about object \emph{ownership} so that memory is released at a
predictable moment without a garbage collector.

\section{Arrays in memory: row-major and column-major}
\label{sec:rowcolumn}

A matrix is two-dimensional, but memory is one-dimensional: a long row of numbered bytes.
There are two ways to lay out a matrix
$A = \begin{pmatrix} a_{11} & a_{12} & a_{13} \\ a_{21} & a_{22} & a_{23}\end{pmatrix}$
in memory:
\begin{center}
\begin{tikzpicture}[cell/.style={draw,minimum width=10mm,minimum height=6mm,font=\small}]
\node[font=\sffamily\small,anchor=east] at (-0.3,0) {row-major (C, C++, \py/NumPy):};
\foreach \v [count=\i] in {$a_{11}$,$a_{12}$,$a_{13}$,$a_{21}$,$a_{22}$,$a_{23}$}
  \node[cell,fill=nublue!10] at (\i*1.0,0) {\v};
\node[font=\sffamily\small,anchor=east] at (-0.3,-1) {column-major (Fortran, MATLAB, R, Julia):};
\foreach \v [count=\i] in {$a_{11}$,$a_{21}$,$a_{12}$,$a_{22}$,$a_{13}$,$a_{23}$}
  \node[cell,fill=nugold!15] at (\i*1.0,-1) {\v};
\end{tikzpicture}
\end{center}
In \term{row-major} order the elements of each row are next to each other; in
\term{column-major} order the elements of each column are. This matters for speed,
because the processor loads memory in chunks (\emph{cache lines} of 64 bytes). A loop that
walks through memory in order is fast; a loop that jumps around is slow. In C you should
make the \emph{last} index change fastest; in Fortran, the \emph{first}.
\begin{lstlisting}[language=C]
/* C: fast, j (last index) varies fastest */
for (i = 0; i < n; i++)
    for (j = 0; j < n; j++)
        s += a[i][j];
\end{lstlisting}
\begin{lstlisting}[language={[08]Fortran}]
! Fortran: fast, i (first index) varies fastest
do j = 1, n
    do i = 1, n
        s = s + a(i, j)
    end do
end do
\end{lstlisting}

\section{Indexing: from 0 or from 1?}

Languages also disagree about the number of the first element. In C, C++, C\#, Java,
JavaScript, \py, Go and Rust the first element of an array is element \code{0}. In
Fortran, MATLAB, R, Julia, the Wolfram Language and Lua it is element \code{1}. (Fortran
even lets you choose any starting index.)

Zero-based indexing comes from C, where \code{a[i]} means ``the element $i$ positions
after the start of the array''. One-based indexing matches mathematical notation, where
the first component of a vector is $x_1$. Neither is better. But when you translate an
algorithm from a textbook (one-based) to \py (zero-based) or from MATLAB to \py, you must
shift every index, and \term{off-by-one errors} are among the most common bugs in
programming.

\begin{pitfall}[Common pitfall: ranges]
\py's \code{range(1, 5)} gives $1, 2, 3, 4$: the end is \emph{excluded}. MATLAB's
\code{1:5} and Julia's \code{1:5} give $1, 2, 3, 4, 5$: the end is \emph{included}.
Fortran's \code{do i = 1, 5} also includes 5. R's \code{1:5} includes 5 too. Check the
first and last iteration of every loop you translate.
\end{pitfall}

\section{Three kinds of errors}

Whatever language you use, you will meet three kinds of errors.
\begin{enumerate}
\item \textbf{Syntax errors}: the code does not follow the grammar of the language
  (a missing bracket, a misspelt keyword). The compiler or interpreter tells you where;
  read the message carefully, and look at the line \emph{before} the one reported too.
\item \textbf{Run-time errors}: the code is grammatically correct but fails while
  running (dividing by zero in integer arithmetic, opening a file that does not exist,
  using an index outside an array). \py raises an \term{exception} with a
  \term{traceback} that shows where it happened.
\item \textbf{Logic errors}: the program runs and produces an answer, but the answer is
  wrong (a wrong sign, a wrong unit, a loop that stops one step early). These are the most
  dangerous because nothing tells you about them. The only defence is \emph{testing}:
  compare with a hand calculation, an exact solution, a limiting case, or a different
  method.
\end{enumerate}

\begin{history}[A little history: the Mars Climate Orbiter]
In September 1999 NASA's Mars Climate Orbiter was lost when it entered the Martian
atmosphere far too low. One team's software reported thruster impulse in pound-force
seconds; the navigation software expected newton-seconds. The program ran without any
error message: it was a logic error, a units mismatch. Some modern tools can attach
physical units to numbers to catch such mistakes (the \py package \code{pint}, units in
the Wolfram Language and F\#, \code{Unitful.jl} in Julia).
\end{history}

\section{Parallel computing}

Since about 2005 the clock speed of processors has hardly increased. Instead, computers
have more and more \emph{cores}: a laptop has 8--16, a workstation 32--128, and a
supercomputer millions. A program that uses only one core uses only a small fraction of the
machine. \term[parallel computing]{Parallel computing} divides the work among many cores, or among many
computers. There are three common models, and you will meet their names in the
documentation of every simulation package.
\begin{itemize}
\item \textbf{Shared memory (threads).} Several threads of one program run on the cores of
  one computer and share its memory. In C, C++ and Fortran the usual tool is
  \textbf{OpenMP}: a comment-like directive before a loop asks the compiler to split the
  iterations among threads. Julia has \code{Threads.@threads}; Java, C\# and Rust have
  thread libraries; NumPy and MATLAB use threads automatically inside many library routines
  (for example, matrix multiplication).
\item \textbf{Distributed memory (message passing).} Many separate processes, possibly on
  different computers of a cluster, each with its own memory, exchange data by sending
  messages. The standard is \textbf{MPI} (Message Passing Interface), used by almost every
  large CFD, FEM, MD and climate code, with bindings for C, C++, Fortran, \py
  (\code{mpi4py}) and Julia.
\item \textbf{Accelerators (GPUs).} A GPU applies the same operation to thousands of data
  elements at once. It is programmed with \textbf{CUDA} (NVIDIA), HIP (AMD), SYCL or OpenCL,
  or from high-level languages through libraries such as CuPy, PyTorch, JAX and
  \code{CUDA.jl}. Molecular dynamics, CFD and machine learning gain factors of ten to a
  hundred.
\end{itemize}
Parallel programs are harder to write and to debug than serial ones. Two threads that
update the same variable at the same time cause a \term{race condition}, whose result
depends on timing and may differ from run to run. And not everything can be parallelised:
if 10\% of a program is inherently serial, then even with infinitely many cores it cannot
run more than 10 times faster (\term{Amdahl's law}). The good news for most engineers is
that parallelism is usually already built into the libraries and simulation tools they use.

\section{Libraries: standing on the shoulders of Fortran}

No engineer writes a matrix inversion or a fast Fourier transform from scratch for real
work. They use \term{libraries}: collections of well-tested routines. A few libraries are
so important that almost every language uses them underneath.
\begin{itemize}
\item \textbf{BLAS} (Basic Linear Algebra Subprograms, 1979 onwards) defines vector and
  matrix operations such as dot products and matrix multiplication.
\item \textbf{LAPACK} (1992) solves linear systems, eigenvalue problems and least-squares
  problems using BLAS. It was originally written in Fortran 77 and is now in modern
  Fortran.
\item Optimised implementations such as \textbf{OpenBLAS}, Intel's \textbf{oneMKL} and
  Apple's \textbf{Accelerate} make these routines run close to the peak speed of the
  hardware.
\item \textbf{FFTW} computes fast Fourier transforms.
\end{itemize}
When you write \code{numpy.linalg.solve(A, b)} in \py, \code{A\textbackslash b} in MATLAB,
\code{solve(A, b)} in R or \code{A \textbackslash\ b} in Julia, the actual work is done by
LAPACK and BLAS. In a sense, every scientific language is a friendly front end to the same
Fortran and C code.

\begin{exercises}
\item Convert the decimal numbers 5, 12 and 100 to binary by hand. Check your answers with
  \py's \code{bin()} function.
\item What is the largest integer that can be stored in an unsigned 16-bit integer? A
  sensor reports a count every second into a 16-bit unsigned counter. After how long does
  the counter overflow?
\item Run \code{0.1 + 0.2 == 0.3} in \py. Then try the same in a web browser's JavaScript
  console (press F12) with \code{0.1 + 0.2 === 0.3}. Explain the result.
\item The quadratic formula $x = (-b \pm\sqrt{b^2-4ac})/(2a)$ suffers from cancellation when
  $b^2 \gg 4ac$. Compute the small root of $x^2 + 10^8 x + 1 = 0$ with the formula in \py
  and compare with the rewritten formula $x = 2c/(-b - \sqrt{b^2-4ac})$. Which is right?
\item Write a \py loop that sums $1/n^2$ for $n = 1, 2, \ldots, 10^6$, first from $n=1$
  upwards and then from $n=10^6$ downwards, using \code{numpy.float32} for the running
  total. Compare both with $\pi^2/6$. Why does the order matter?
\item A \py student writes \code{a = [[0]*3]*3} to create a $3\times3$ matrix of zeros
  and then sets \code{a[0][0] = 1}. Print \code{a}. Explain the surprising result using the
  idea of references. How should the matrix be created?
\end{exercises}
